\documentclass[11pt]{article}
\usepackage{array}
\usepackage{amsmath}
\usepackage{booktabs}
\usepackage{geometry}
\geometry{margin=1in}

\title{STAT 17 Discussion Notes \\ Week 10}
\date{March 2026}
\author{Xiao Zhang}

\begin{document}

\maketitle

\section{Comparing Two Groups}

When comparing two groups, the method depends on the type of response variable and whether the samples are independent or paired.

\begin{center}
\begin{tabular}{c|c|c|c}
\toprule
 & Two Proportions & Two Means (Independent Samples) & Two Means (Dependent) \\
\midrule
Parameter & $p_1 - p_2$ & $\mu_1 - \mu_2$ & $\mu_d$ \\

Response Variable & Qualitative & Quantitative & Quantitative \\

Sampling Structure & Independent samples & Independent samples & Matched / paired data \\

Distribution & Standard Normal ($z$) & $t$ distribution & $t$ distribution ($df=n-1$) \\

\bottomrule
\end{tabular}
\end{center}


\section{Two-Proportion Inference}

\subsection*{Conditions}

\begin{itemize}
\item Independent random samples
\item Sample size less than $5\%$ of population
\item Success-failure condition
\end{itemize}

\[
n_1\hat p_1(1-\hat p_1) \ge 10
\]

\[
n_2\hat p_2(1-\hat p_2) \ge 10
\]

\subsection*{Confidence Interval}

\[
(\hat p_1 - \hat p_2)
\pm
z_{\alpha/2}
\sqrt{
\frac{\hat p_1(1-\hat p_1)}{n_1}
+
\frac{\hat p_2(1-\hat p_2)}{n_2}
}
\]

\subsection*{Test Statistic}

\[
z =
\frac{\hat p_1 - \hat p_2}
{\sqrt{\hat p(1-\hat p)
\left(\frac{1}{n_1}+\frac{1}{n_2}\right)}}
\]

where

\[
\hat p = \frac{x_1+x_2}{n_1+n_2}
\]

\section{Two-Sample $t$ Procedures}

\subsection*{Conditions}

\begin{itemize}
\item Independent random samples
\item Sample size less than $5\%$ of population
\item Nearly normal populations OR large samples ($n \ge 30$)
\end{itemize}

\subsection*{Confidence Interval}

\[
(\bar{x}_1 - \bar{x}_2)
\pm
t_{\alpha/2}
\sqrt{\frac{s_1^2}{n_1} + \frac{s_2^2}{n_2}}
\]

\subsection*{Test Statistic}

\[
t =
\frac{\bar{x}_1 - \bar{x}_2}
{\sqrt{\frac{s_1^2}{n_1} + \frac{s_2^2}{n_2}}}
\]

\section{Paired $t$ Procedures}

Used when observations are paired (e.g., before–after measurements or matched subjects).

Let

\[
d = x_1 - x_2
\]

where

\begin{itemize}
\item $d$ = difference for each pair
\item $x_1$ = first observation in the pair
\item $x_2$ = second observation in the pair
\end{itemize}

We then treat the differences as a single sample.

\subsection*{Notation}

\begin{itemize}
\item $\bar d$ = sample mean of the differences
\item $s_d$ = sample standard deviation of the differences
\item $n$ = number of pairs
\end{itemize}

\subsection*{Confidence Interval}

\[
\bar d \pm t_{\alpha/2}\frac{s_d}{\sqrt{n}}
\]

\subsection*{Test Statistic}

\[
t =
\frac{\bar d - 0}{s_d/\sqrt{n}}
\]

\section{Chi-Square Test of Independence}

Used to determine whether two categorical variables are associated.

\subsection*{Hypotheses}

\[
H_0: \text{The two variables are independent}
\]

\[
H_1: \text{The two variables are dependent}
\]

\subsection*{Conditions}

\begin{itemize}
\item All expected counts $\ge 1$
\item No more than 20\% of expected counts $< 5$
\end{itemize}

\subsection*{Expected Counts}

\[
E = \frac{(\text{row total})(\text{column total})}{n}
\]

\subsection*{Test Statistic}

\[
\chi^2 =
\sum \frac{(O - E)^2}{E}
\]

\subsection*{Degrees of Freedom}

\[
df = (r-1)(c-1)
\]

\subsection*{P-value}

\[
p = P(\chi^2 > \chi_0^2)
\]

\section{Linear Regression}

Linear regression describes the relationship between two quantitative variables.

\begin{itemize}
\item Explanatory variable: $x$
\item Response variable: $y$
\end{itemize}

\subsection*{Regression Line}

\[
y = mx + b
\]

where

\begin{itemize}
\item $m$ = slope
\item $b$ = $y$-intercept
\end{itemize}

For prediction we use

\[
\hat{y} = mx + b
\]

where $\hat{y}$ is the predicted value of $y$.

\subsection*{Interpretation of the Slope}

The slope $m$ represents the expected change in the response variable $y$ for a one-unit increase in the explanatory variable $x$.

In context:

\begin{quote}
For each one-unit increase in $x$, the predicted value of $y$ changes by $m$ units on average.
\end{quote}

\subsection*{Interpretation of the Intercept}

The intercept $b$ represents the predicted value of $y$ when $x = 0$.

\begin{quote}
When $x = 0$, the predicted value of the response variable is $b$.
\end{quote}

Note: Sometimes $x=0$ may not be meaningful in the context of the data.

\subsection*{Residual}

A residual measures the difference between the observed value and the predicted value.

\[
\text{Residual} = y - \hat{y}
\]

where

\begin{itemize}
\item $y$ = observed value
\item $\hat{y}$ = predicted value from the regression line
\end{itemize}

Residuals are used to evaluate how well the regression line fits the data.


\end{document}
